% ============================================================================
%  Chapitre 2, partie III : Contraintes
%  (manuscrit P2, p. 4 à 10)
% ============================================================================

% ----------------------------------------------------------------------------
\section{Contraintes}
% ----------------------------------------------------------------------------

% ----------------------------------------------------------------------------
\subsection{Configuration et mouvement admissible}

\begin{definition}[Configuration]
    Une \emph{configuration} décrit la position géométrique complète du
    système à un instant donné. L'ensemble des configurations constitue
    l'\emph{espace de configuration}, noté $Q$.
\end{definition}

Pour un solide libre dans l'espace,
\[
    Q = \underbrace{\R^3}_{\text{position d'un point fixe du solide}}
    \times
    \underbrace{\SO}_{\text{orientation}}
\]
soit $3 + 3 = 6$ degrés de liberté. On utilise $q$ pour désigner une
configuration ou ses coordonnées généralisées, à ne pas confondre avec
le quaternion $\quat{q}$.

\begin{definition}[Coordonnées généralisées]
    Les \emph{coordonnées généralisées} sont un ensemble de variables
    indépendantes qui décrivent entièrement la configuration d'un
    système physique en tenant compte de ses contraintes.
\end{definition}

Exemple : Pour un pendule double plan, les deux masses ont pour positions
$\vect{r}_1 = (x_1,y_1)$ et $\vect{r}_2 = (x_2,y_2)$, soit $4$
coordonnées. Les deux longueurs fixées imposent $2$ contraintes, et il
suffit de $4 - 2 = 2$ coordonnées généralisées : $(q_1,q_2) =
(\theta_1,\theta_2)$.

\begin{definition}[Mouvement]
    Un \emph{mouvement} est une courbe $t\mapsto q(t)\in Q$. Sa vitesse
    appartient à l'espace tangent à $Q$ :
    \[
        \dot q(t)\in T_{q(t)}Q
    \]
\end{definition}

Pour un solide, cette vitesse est représentée par $(\vect{v}_O,
\vect{\Omega})$. Une contrainte restreint les configurations
possibles, les vitesses possibles, ou les deux. Un mouvement qui
respecte ces restrictions est dit \emph{cinématiquement admissible}.

% ----------------------------------------------------------------------------
\subsection{Contraintes holonomes}

\paragraph{Définition}

\begin{definition}[Contrainte holonome]
    Une contrainte est \emph{holonome} lorsqu'elle s'exprime par une
    relation entre les coordonnées et éventuellement le temps. Pour
    $k$ contraintes,
    \[
        f(q,t) = \begin{pmatrix}
            f_1(q,t)\\ \vdots\\ f_k(q,t)
        \end{pmatrix} = 0,
        \qquad f : \R^n\times\R \to \R^k
    \]
    Elles sont \emph{indépendantes} si $\rang D_q f = k$, avec
    $D_q f\in\mathcal{M}_{k,n}(\R)$.
\end{definition}

Avec $q\in\R^n$, soit $n$ degrés de liberté, et $k$ contraintes
indépendantes $f(q) = 0$, l'ensemble des configurations admissibles
est
\[
    \mathcal{C} = \{q\in\R^n : f(q) = 0\}
\]
D'après le théorème du rang appliqué à $D_q f$,
\[
    \underbrace{\rang D_q f}_{k\ \text{contraintes}}
    + \underbrace{\dim\Ker D_q f}_{\text{directions admissibles}}
    = \underbrace{n}_{\text{degrés de liberté}}
\]
Il suffit donc de $n - k$ coordonnées généralisées pour décrire la
cinématique admissible.

\paragraph{Exemple : pendule plan}

Les coordonnées $(x,y)$ de la masse vérifient $x^2 + y^2 = \ell^2$ ;
la contrainte est holonome,
\[
    f(x,y) = x^2 + y^2 - \ell^2 = 0
\]
et l'angle $\theta$ suffit à décrire le système. Si $q(t)$ respecte
$f(q(t),t) = 0$, en dérivant,
\[
    D_q f(q,t)\,\dot q + \pd{f}{t}(q,t) = 0
\]
Ici $f$ ne dépend pas du temps, donc
\[
    D_q f(q)\,\dot q = 0
    \longleftrightarrow
    \begin{pmatrix} 2x & 2y \end{pmatrix}
    \begin{pmatrix} \dot x\\ \dot y \end{pmatrix} = 0
    \longleftrightarrow
    2\,\langle (x,y),(\dot x,\dot y)\rangle = 0
\]
La vitesse est perpendiculaire à la corde.

% ----------------------------------------------------------------------------
\subsection{Contraintes sur les vitesses et non-holonomie}

\paragraph{Forme générale}

De nombreuses liaisons mécaniques s'écrivent sous forme de vitesses :
\[
    \mat{A}(q,t)\,\dot q + b(q,t) = 0
\]
Comme on n'étudie que la sphère sur un sol fixe, on ne considérera
que la forme
\[
    \mat{A}(q)\,\dot q = 0,
    \qquad \mat{A}(q)\in\mathcal{M}_{k,n}(\R)
\]

\paragraph{Non-holonomie}

\begin{definition}[Contrainte non holonome]
    Une contrainte sur les vitesses est \emph{non holonome} lorsqu'elle
    ne peut pas être remplacée, même localement, par des contraintes
    équivalentes portant uniquement sur les configurations.
\end{definition}

Si $\mat{A}(q)$ est de rang $k$, les vitesses admissibles
$\Ker\mat{A}(q)$ forment un espace de dimension $n - k$. Pour une
contrainte non holonome, on ne peut donc pas trouver $k$ fonctions
indépendantes $f_i$ telles que
\[
    \Ker\mat{A}(q) = \Ker D f(q)
\]

% ----------------------------------------------------------------------------
\subsection{Déplacements virtuels et réactions idéales}

\paragraph{Déplacement virtuel}

\begin{definition}[Déplacement virtuel]
    Un \emph{déplacement virtuel} $\delta q$ est une variation
    infinitésimale de la configuration à temps fixé, $\delta t = 0$.
\end{definition}

Pour une contrainte holonome $f(q,t) = 0$, les déplacements virtuels
admissibles vérifient
\[
    D_q f(q,t)\,\delta q = 0
\]
Pour une contrainte cinématique $\mat{A}(q)\,\dot q = 0$ (roulement,
glissement), on fige le temps et on impose de même
\[
    \mat{A}(q)\,\delta q = 0
\]

\paragraph{Réaction idéale}

On note $\vect{Q}^r$ les forces généralisées de réaction. Leur travail virtuel
est
\[
    \delta W_r = (\vect{Q}^r)^T\delta q
\]

\begin{definition}[Liaison idéale]
    Une liaison est \emph{idéale} lorsque ses réactions ont un travail
    virtuel nul pour tout déplacement virtuel admissible :
    \[
        (\vect{Q}^r)^T\delta q = 0
        \qquad \forall\,\delta q\in\Ker\mat{A}(q)
    \]
\end{definition}

\begin{proposition}[Multiplicateurs de Lagrange]
    Pour une liaison idéale, il existe $\vect{\lambda}\in\R^k$ tel que
    \[
        \vect{Q}^r = \mat{A}(q)^T\vect{\lambda}
    \]
\end{proposition}

\begin{demonstration}
    $\vect{Q}^r\in(\Ker\mat{A}(q))^\perp = \Ima\mat{A}(q)^T$ par
    l'orthogonalité noyau-image.
\end{demonstration}

\begin{note}
    $\mat{A}(q)$ envoie les vitesses généralisées ($\R^n$) sur les
    vitesses contraintes ($\R^k$). À l'inverse, $\mat{A}(q)^T$ envoie
    les forces de contrainte $\vect{\lambda}\in\R^k$ sur les forces
    généralisées ($\R^n$) : à partir des forces de contrainte, on
    reconstruit les forces généralisées.
\end{note}

\newpage
